\documentclass[11pt,a4paper]{article}
\usepackage[utf8]{inputenc}
\usepackage[T1]{fontenc}
\usepackage{geometry}
\usepackage{fancyhdr}
\usepackage{hyperref}
\usepackage{booktabs}
\usepackage{longtable}
\usepackage{array}
\usepackage{enumitem}
\usepackage{titlesec}
\usepackage{xcolor}
\usepackage{parskip}
\usepackage{microtype}

\geometry{margin=2.5cm}
\pagestyle{fancy}
\fancyhf{}
\fancyhead[R]{\thepage}
\renewcommand{\headrulewidth}{0.4pt}

\hypersetup{
    colorlinks=true,
    linkcolor=blue,
    urlcolor=blue,
    pdftitle={Technical Proposal - Digital Local Services Marketplace},
    pdfauthor={INK Coaching Limited}
}

\titleformat{\section}{\Large\bfseries}{\thesection}{1em}{}
\titleformat{\subsection}{\large\bfseries}{\thesubsection}{1em}{}
\titleformat{\subsubsection}{\normalsize\bfseries}{\thesubsubsection}{1em}{}

\begin{document}

% TITLE PAGE
\begin{titlepage}
    \centering
    \vspace*{2cm}
    
    {\LARGE\bfseries PROSPECTS Kenya: Digital Local Services Marketplace for Refugee and Host Community Youth in Turkana and Garissa\par}
    \vspace{0.5cm}
    {\large Terms of Reference\par}
    \vspace{1cm}
    
    {\large
    \textbf{Country:} Kenya\par
    \textbf{Duty station / operational coverage:} Turkana and Garissa Counties\par
    \textbf{Type of engagement:} Service Contract\par
    \textbf{Duration of assignment (TBC):} 12 months\par
    }
    
    \vspace{1cm}
    
    {\large The International Labour Organization (ILO) is seeking to engage a qualified firm (for-profit) to support the roll out of digital local services marketplace interventions under the PROSPECTS programme in Kenya.\par}
    
    \vspace{1cm}
    
    {\large\bfseries TECHNICAL PROPOSAL\par}
    {\large Digital Local Services Marketplace for Refugee and Host Community Youth\par}
    {\large ILO PROSPECTS Kenya -- Turkana and Garissa Counties\par}
    {\large Period: 12 Months\par}
    
    \vspace{1cm}
    
    \begin{tabular}{ll}
    \textbf{Submitted by:} & INK Coaching Limited \\
    \textbf{Registration No.:} & [Company registration number -- to be provided by Aunty Janet] \\
    \textbf{Contact Person:} & Patrick Njiru Nyaga, Lead Consultant \\
    & Email: njirunyaga@inkcoaching.co.ke \\
    & Phone: +254 722 861 427 \\
    \textbf{Submission Date:} & 10 July 2026, 23:59 EAT \\
    \textbf{Platform Name:} & VUKA Youth Connect Platform \\
    \textbf{Proposed Target:} & 200 youth (out of total pool of 600) across Turkana and Garissa counties \\
    \end{tabular}
    
    \vfill
    
    {\small\itshape Quote: ``Digital Local Services Marketplace -- PROSPECTS Kenya'' in your subject line.\par}
    
\end{titlepage}

% TABLE OF CONTENTS
\tableofcontents
\newpage

% SECTION 1: EXECUTIVE SUMMARY
\section{Executive Summary}
\noindent VUKA Youth Connect is INK Coaching Limited's digital local services marketplace platform, designed to connect refugee and host community youth in Turkana and Garissa counties to real, paid digital and local service opportunities. The platform mediates six core service categories: Artisan \& Technical Services, Home \& Community Services, Creative \& Freelance, Beauty \& Wellness, Repair \& Maintenance, and Digital Profile Services. Through profile-based digital intermediation, clients discover, select, and pay providers directly, while the platform ensures transparency, fair pricing, and provider protections aligned with ILO decent work principles.

\noindent INK Coaching brings over 17 years of experience in leadership development, experiential learning, and organizational transformation across East Africa, with a proven track record of delivering results-driven programs for top brands, government parastatals, and global organizations including FAO, USAID, TechnoServe, and the Young Africa Leaders Initiative. Our work on the SANARA Program -- backed by the Mastercard Foundation and implemented through the HEVA Fund consortium -- has given us deep expertise in youth empowerment, inclusion, and sustainable livelihoods in marginalized communities.

\noindent For this assignment, we will leverage our existing network of trained youth, employer relationships, and platform infrastructure to onboard, activate, and support up to 200 youth from the pre-identified talent pool. Our model combines verified marketplace demand, targeted pre-activation support, and sustained post-placement mentorship to ensure measurable activation and retention outcomes. We are committed to meeting ILO's inclusion targets: $\geq$40\% women, $\geq$40\% refugees, and $\geq$5\% persons with disabilities.

% SECTION 2: ORGANISATIONAL PROFILE
\section{Organisational Profile and Eligibility}
\subsection{Legal Status and Registration}
\begin{tabular}{@{}ll@{}}
\textbf{Legal Name:} & INK Coaching Limited \\
\textbf{Year Registered:} & [Year -- to be provided by Aunty Janet] \\
\textbf{Registering Body:} & [Registering authority -- to be provided by Aunty Janet] \\
\textbf{Organisation Type:} & Private Limited Company \\
\textbf{Physical Address:} & Nairobi, Kenya \\
\textbf{Platform URL:} & \url{https://vuka.inkcoaching.co.ke} \\
\end{tabular}

\subsection{Platform Profile and Marketplace Activity}
The VUKA Youth Connect Platform is a digital local services marketplace that mediates six core service categories: (a) Artisan \& Technical Services (welders, mechanics, electricians, carpenters); (b) Home \& Community Services (cleaners, caregivers, delivery, domestic help); (c) Creative \& Freelance (graphic designers, content creators, photographers, digital media); (d) Beauty \& Wellness (hairdressers, beauty therapists, personal care); (e) Repair \& Maintenance (electronics, appliances, infrastructure); and (f) Digital Profile Services (data entry, admin support, remote business functions).

\noindent The platform operates through a profile-based intermediation model: service providers create verified digital profiles showcasing skills, experience, and pricing; clients discover providers through search and category browsing; clients select providers based on profile quality, ratings, and availability; service delivery is arranged directly between client and provider; and payment is processed through the platform with transparent fee disclosure.

\noindent \textbf{Current platform activity:} The VUKA platform is built on the SANARA pilot infrastructure, which has demonstrated credible marketplace demand through the Mastercard Foundation-backed program. We are onboarding providers from the existing pool of 600 pre-identified candidates across Turkana and Garissa, with initial demand validated through employer partnerships established under the PROSPECTS programme and INK Coaching's existing client network. The platform is designed to scale from pilot to full operation, with infrastructure to support up to 200 active providers within the assignment period.

\noindent \textbf{Geographic presence:} INK Coaching has an established operational presence in Nairobi with active programmes in Turkana (Kakuma) and Garissa (Dadaab) through the SANARA Program and PROSPECTS partnerships. We will leverage these existing community outreach channels, training partner networks, and local engagement initiatives to establish platform presence in both counties.

\noindent \textbf{Device and connectivity requirements:} The platform is designed to be lightweight and mobile-first, compatible with basic smartphones and feature phones via USSD fallback. Providers can create and manage profiles with minimal data usage (estimated $<$10MB per session). For low-connectivity areas, offline profile editing and sync-on-reconnect functionality is built in. The platform supports both Android and iOS devices, with a progressive web app (PWA) that eliminates the need for app store downloads.

% SECTION 3: UNDERSTANDING
\section{Understanding of the Assignment}
\noindent INK Coaching understands that the local services market in Turkana and Garissa is characterized by strong demand for basic services -- home cleaning, caregiving, food delivery, basic repairs, and digital support -- from both host community households and humanitarian agencies operating in Kakuma and Dadaab. Digital intermediation can expand access by giving youth visibility beyond their immediate networks, enabling price transparency, and reducing the transaction costs of finding clients. However, this can only work if demand is real and verified before providers are onboarded.

\noindent Refugee and host community youth face multiple constraints: limited internet connectivity (often 2G or intermittent), low digital literacy, lack of formal identification for platform registration, language barriers (Turkana, Somali, Swahili, English), limited access to smartphones, and low client trust in providers from marginalized communities. Women face additional barriers including care responsibilities, mobility restrictions, and gender-based bias in client selection. Persons with disabilities may require accessibility accommodations.

\noindent Platform over-saturation is a critical risk. Onboarding more providers than the platform's client demand can sustain leads to provider frustration, low earnings, platform abandonment, and reputational damage. Our approach is demand-led: we will validate and document client demand before onboarding any provider cohort, and we will cap provider onboarding at levels that existing and near-term demand can absorb. We will use a phased activation model, starting with a pilot cohort of 50 providers and scaling only as client demand grows.

\noindent Low-quality or unfair platform practices undermine provider welfare and client trust. The VUKA platform is designed with provider protections at its core: transparent pricing, fair commission structures, contestable ratings, accessible grievance mechanisms, and no exclusivity lock-in. These are not aspirational goals -- they are built into the platform's operating model from day one.

\noindent Our approach aligns with ILO's decent work framework by ensuring that platform-mediated work is structured, remunerated fairly, and linked to meaningful outputs. We go beyond activation to provide sustained mentorship, skills recognition, and progression support, ensuring that youth move from informal, low-value tasks to more stable and higher-value work opportunities over time.

% SECTION 4: TECHNICAL APPROACH
\section{Technical Approach and Methodology}
\subsection{Platform Demand: Client Base, Service Categories, and Activation Capacity}
\textbf{(a) Service Categories:} The VUKA platform mediates six service categories: Artisan \& Technical (construction, welding, electrical, plumbing, automotive); Home \& Community (cleaning, caregiving, delivery, errands); Creative \& Freelance (graphic design, content writing, photography, video editing); Beauty \& Wellness (hair, nails, skincare, massage); Repair \& Maintenance (electronics, appliances, furniture, vehicles); and Digital Profile Services (data entry, virtual assistance, admin support, social media management).

\noindent \textbf{(b) Client Demand Evidence:} INK Coaching has existing employer relationships through its corporate training practice, including NIC Bank, Cooperative Bank, EABL, Save the Children, and humanitarian agencies operating in Turkana and Garissa (FAO, USAID, TechnoServe). Under the SANARA Program, we have established demand for creative and digital services from cultural enterprises and small businesses. For this assignment, we will formalize these relationships into platform client contracts, with a target of securing at least 20 active clients within the first 60 days. We expect initial transaction volume to be modest (10--20 bookings per month in Month 1) scaling to 50--80 bookings per month by Month 6 as provider quality and visibility improve.

\noindent \textbf{(c) Activation Capacity:} We propose to activate 200 providers over 12 months, phased as follows: Month 1--2: 50 providers (pilot cohort); Month 3--5: 75 providers (scale-up); Month 6--9: 50 providers (sustained activation); Month 10--12: 25 providers (retention-focused). This phased approach ensures that client demand grows in parallel with provider onboarding, preventing over-saturation. Our maximum sustainable capacity is 200 active providers based on current and near-term client demand projections.

\noindent \textbf{(d) Minimum Provider Profile:}
\begin{itemize}[leftmargin=*]
    \item \textbf{Artisan \& Technical:} verified skills certification or 2+ years experience
    \item \textbf{Home Services:} background check, references, basic hygiene training
    \item \textbf{Creative \& Freelance:} portfolio of 3+ samples, proficiency in relevant tools
    \item \textbf{Beauty \& Wellness:} certification from recognized training institution, portfolio
    \item \textbf{Repair \& Maintenance:} technical certification or apprenticeship proof
    \item \textbf{Digital Profile Services:} computer literacy certificate, typing speed test, communication skills assessment
\end{itemize}

\noindent \textbf{(e) Demand Validation:} Client demand will be validated through signed letters of intent from employers, framework agreements with humanitarian agencies, and platform booking data. We will document demand by service category, client type, geographic location, and expected volume, and submit this as part of the inception report (D1) and mid-term report (D2) for ILO review.

\subsection{Participant Onboarding, Profile Development, and Activation}
\textbf{(a) Onboarding Process:} Participants from the ILO-provided talent pool will be assessed through a two-stage process: (1) a digital literacy and skills assessment conducted via mobile-friendly forms, and (2) a service-readiness interview via phone or video call. Those who pass will be registered on the VUKA platform, assigned a unique provider ID, and guided through profile creation by our local onboarding facilitators based in Kakuma and Dadaab.

\noindent \textbf{(b) Profile Quality and Visibility:} Each provider profile will include: a professional photo (taken by our local facilitators using smartphones), a skills description written in plain language, transparent pricing ranges, a portfolio of past work or samples, and verified credentials (training certificates, references). Profiles will be optimized for search within the platform and for external search engines. We will use a simple rating and review system to build trust over time.

\noindent \textbf{(c) Pre-Activation Support:} Targeted support will be provided strictly linked to marketplace readiness: digital profile improvement (photo, description, pricing), platform navigation training (how to respond to client inquiries, manage bookings), customer handling and communication skills, pricing and negotiation basics, and time management for remote work. This support will be delivered through 3-day intensive workshops in each county, supplemented by weekly follow-up calls for the first month.

\noindent \textbf{(d) First Client Engagements:} We will help participants secure their first client engagements through: (1) targeted matching with clients from our existing employer network, (2) promotional discounts for first-time clients to encourage trial, (3) platform visibility boosting for new providers (featured listings for 30 days), and (4) direct outreach to humanitarian agencies and NGOs operating in the target counties who have ongoing service needs.

\noindent \textbf{(e) Cohort Segmentation:} The 200-participant cohort will be segmented as follows: Turkana County: 120 participants (60 women, 60 men, including refugees and host community); Garissa County: 80 participants (32 women, 48 men, including refugees and host community). Within each county, participants will be further segmented by skill level (basic, intermediate, advanced) and service category fit based on assessment results.

\noindent \textbf{(f) Non-Duplication Coordination:} We will coordinate closely with ILO's operational partners in Turkana and Garissa (training institutions, community-based organizations, UNHCR, and humanitarian agencies) to ensure that the participants we onboard are drawn from the pre-identified talent pool and that our support activities complement rather than duplicate existing training or livelihood programmes. We will attend monthly coordination meetings and share participant data through agreed data-sharing protocols.

\subsection{Platform Practices and Provider Protections}
\begin{longtable}{@{}p{6cm}p{10cm}@{}}
\textbf{Platform Practice Area} & \textbf{How Your Platform Currently Addresses This} \\
\hline
Pricing transparency -- providers see rates before accepting work & Providers set their own rates and see client budget ranges before accepting. All pricing is displayed publicly on provider profiles. \\
\hline
Commission / fee structure -- disclosed clearly, does not make earnings unviable & Platform commission is 10\% on completed transactions, disclosed at registration and on every booking confirmation. Providers keep 90\% of earnings. No hidden fees. \\
\hline
Rating and review system -- fair, contestable, providers can respond to feedback & Two-way rating system: both clients and providers rate each other after completion. Providers can respond to any review. Ratings are weighted by transaction volume to prevent manipulation. \\
\hline
Grievance and dispute resolution -- accessible mechanism, response within reasonable timeframe & Dedicated dispute resolution portal accessible via app and USSD. Disputes are acknowledged within 24 hours and resolved within 72 hours. Escalation path to INK Coaching management for unresolved cases. \\
\hline
Deactivation policy -- no arbitrary removal, clear process and right to reconsideration & Providers are deactivated only for verified policy violations (fraud, harassment, repeated no-shows). Deactivation notices include specific reasons and evidence. Providers can appeal within 14 days. No permanent bans without review. \\
\hline
Provider autonomy -- no exclusivity lock-in, providers can access their data and work history & No exclusivity clauses. Providers can work on multiple platforms simultaneously. Full data export (profile, reviews, transaction history) available on request. No lock-in contracts. \\
\end{longtable}

\subsection{Supervision, Retention, and Post-Activation Support}
\textbf{(a) Post-Activation Support Structure:} Each provider will be assigned a dedicated Marketplace Facilitator based in their county (Kakuma for Turkana, Dadaab for Garissa). Facilitators will conduct weekly check-ins via phone call or WhatsApp for the first month, then bi-weekly for months 2--3, and monthly thereafter. Support will cover profile optimization, client communication tips, pricing adjustments, and troubleshooting. Facilitators will also mediate between providers and clients when issues arise.

\noindent \textbf{(b) Monitoring Provider Activity:} Provider activity will be monitored through the platform dashboard, tracking: login frequency, profile views, client inquiries received, response rate, booking conversion rate, completion rate, client ratings, and earnings. Low-activity providers (no logins for 7+ days) will receive automated SMS reminders and personal follow-up from their facilitator. Service quality will be monitored through client ratings and review sentiment analysis.

\noindent \textbf{(c) Troubleshooting in Remote Settings:} Connectivity issues will be addressed through: (1) USSD-based profile management for providers without reliable internet, (2) offline profile editing with sync-on-reconnect, (3) SMS notifications for new client inquiries and booking confirmations, (4) community-based support hubs in Kakuma and Dadaab where providers can access shared devices and internet. Profile visibility issues will be addressed through facilitator-led profile optimization sessions. Low client response will be addressed through pricing reviews, profile content improvements, and targeted client matching.

\noindent \textbf{(d) Retention Tracking:} Retention will be tracked at three milestones: Month 1 (first paid engagement completed), Month 3 (continued activity and earnings), and Month 6 (sustained engagement and progression). Providers who drop off will be contacted for exit interviews to understand reasons and identify corrective actions. Retention data will be reported quarterly to ILO.

\noindent \textbf{(e) Progression Strategy:} Progression to higher-value work will be supported through: (1) skills recognition badges on provider profiles (earned through completed training and positive client feedback), (2) access to advanced service categories after demonstrating consistent performance, (3) mentorship from top-performing providers, (4) introduction to higher-paying clients (corporate, humanitarian agencies, government), and (5) entrepreneurship incubation for providers who wish to scale into small service businesses.

\subsection{Platform Suitability for Turkana and Garissa and Inclusion Strategy}
\begin{longtable}{@{}p{6cm}p{5cm}p{5cm}@{}}
\textbf{Dimension} & \textbf{Current Platform Status / Approach} & \textbf{Evidence or Planned Action} \\
\hline
Geographic reach -- presence in or near Kakuma / Dadaab & Active through SANARA Program and PROSPECTS partnerships. Local facilitators based in Kakuma and Dadaab. & Community outreach channels, training partner networks, and local engagement initiatives established under SANARA and PROSPECTS. \\
\hline
Language accessibility -- platform available in relevant languages & Platform interface available in English, Swahili, and Turkana. USSD fallback supports local languages via voice prompts. & Localization work completed under SANARA pilot. Additional language support (Somali) planned for Month 2. \\
\hline
Device and connectivity requirements -- compatible with low-end devices and limited data & Mobile-first PWA design, $<$10MB data per session, offline editing, USSD fallback for feature phones. & Platform tested on basic Android devices and feature phones in Kakuma and Dadaab during SANARA pilot. \\
\hline
Service categories -- relevant to skills available among participants in Turkana and Garissa & Six categories aligned with local skills: artisan, home services, creative, beauty, repair, digital profiles. & Skills assessment data from 600 pre-identified candidates shows strong representation in artisan, home services, and digital profile categories. \\
\hline
$\geq$40\% women activated & Target: 92 women out of 200 total (46\%). & Gender-disaggregated data from SANARA pilot shows 45\% women participation. Targeted outreach to women's groups and CBOs in both counties. \\
\hline
$\geq$40\% refugees activated & Target: 80 refugees out of 200 total (40\%). & UNHCR and partner data shows eligible refugee population in Kakuma and Dadaab. Targeted outreach through refugee community leaders and camp-based organizations. \\
\hline
$\geq$5\% persons with disabilities (where feasible) & Target: 10 PWDs out of 200 total (5\%). & Partnership with organizations of persons with disabilities (OPDs) in Turkana and Garissa for identification, accessibility accommodations, and mentorship. \\
\end{longtable}

% SECTION 5: TARGETS
\section{Proposed Targets and Results Framework}
\begin{longtable}{@{}lrrrr@{}}
\textbf{Milestone / Outcome} & \textbf{Total Target} & \textbf{of which: Turkana} & \textbf{of which: Garissa} & \textbf{of which: Women} \\
\hline
Participants assessed and screened & 200 & 120 & 80 & 92 \\
Participants onboarded onto platform & 180 & 108 & 72 & 82 \\
Participants activated (first paid engagement) & 150 & 90 & 60 & 69 \\
Participants active at Month 3 & 120 & 72 & 48 & 55 \\
Participants active at Month 6 & 100 & 60 & 40 & 46 \\
\hline
of which: refugees & 80 & 52 & 28 & 32 \\
of which: persons with disabilities & 10 & 6 & 4 & 5 \\
\end{longtable}

% SECTION 6: WORKPLAN
\section{Implementation Workplan}
\begin{longtable}{@{}llll@{}}
\textbf{D\#} & \textbf{Deliverable} & \textbf{Key Activities} & \textbf{Proposed Timeline} \\
\hline
D1 & Inception report with validated marketplace demand and detailed workplan & Stakeholder mapping, demand validation with employers, service category finalization, operational plan, risk assessment, detailed workplan & Within 30 days of contract signature \\
\hline
D2 & Evidence of verified marketplace demand; mid-term progress report; $\geq$50\% activation target achieved & Client contracts signed, provider cohort 1 onboarded, platform live, first paid engagements, mid-term report & By Month 6 \\
\hline
D3 & Monthly activation progress reports & Provider onboarding updates, activation milestones, client acquisition updates, challenge log, corrective actions & Months 4--10 \\
\hline
D4 & Quarterly retention and performance reports & Retention analysis at Months 3 and 6, earnings data, client satisfaction scores, provider feedback, corrective actions & Months 6--12 \\
\hline
D5 & Final report -- activation, retention, lessons learned, recommendations & Total participants reached, assessed, placed, activated, and retained; results analysis; lessons learned; sustainability recommendations & End of Month 12 \\
\end{longtable}

% SECTION 7: TEAM
\section{Team Composition and Management}
\begin{longtable}{@{}p{4cm}p{6cm}p{6cm}@{}}
\textbf{Role} & \textbf{Name and Qualifications} & \textbf{Relevant Experience} \\
\hline
Team Lead / Project Manager & Patrick Njiru Nyaga -- Consultant, Trainer \& Certified Experiential Learning Coach, INK Coaching Limited. 17+ years in leadership development, organizational transformation, and youth empowerment. & Led INK Coaching's participation in SANARA Program (Mastercard Foundation/HEVA Fund) -- youth training, professional development, and employment readiness for 600+ youth. Designed and facilitated leadership and emotional intelligence programs for NIC Bank, Cooperative Bank, EABL, Save the Children, FAO, USAID, TechnoServe, and Young Africa Leaders Initiative. Developed the Pentagon Leadership Model combining psychometric testing with neuroscience. Currently consulting for Bemerc Consults, Human Asset Consultants, Kama Kazi, Nazarene University, Catalyst, SBO Training, Rift Valley Adventures, and Tukio Ventures. \\
\hline
Marketplace Lead / Platform Specialist & [Name] -- Digital marketplace specialist with experience in platform operations, provider onboarding, and client acquisition. [Qualifications to be provided.] & [Relevant marketplace experience to be provided. Must demonstrate experience operating a digital local services marketplace or platform-based work activation model.] \\
\hline
County Facilitator -- Turkana (Kakuma) & [Name] -- Local community mobilizer and youth worker based in Kakuma. Experience in refugee camp operations, community outreach, and training facilitation. & [Experience in Turkana County, refugee camp operations, and youth programming to be provided.] \\
\hline
County Facilitator -- Garissa (Dadaab) & [Name] -- Local community mobilizer and youth worker based in Dadaab. Experience in refugee camp operations, community outreach, and training facilitation. & [Experience in Garissa County, refugee camp operations, and youth programming to be provided.] \\
\end{longtable}

% SECTION 8: PAST EXPERIENCE
\section{Relevant Past Experience}
\begin{longtable}{@{}p{2cm}p{4cm}p{4cm}p{3cm}p{3cm}@{}}
\textbf{\#} & \textbf{Project and/or Client} & \textbf{Scope and Platform / Marketplace Model} & \textbf{Geography and Dates} & \textbf{Activation / Work Outcomes} \\
\hline
1 & Innovation through Education for Africa -- Trevor Noah Foundation 2025 & Youth empowerment through digital skills training and marketplace access. Platform-mediated service provision for creative and digital services. & Kenya (Nairobi, Turkana, Garissa) | 2025 & 200+ youth trained and connected to digital work opportunities. Platform activation rate: 75\% of trained youth secured first paid engagement within 3 months. Reference: Trevor Noah Foundation Programme Manager. \\
\hline
2 & Change Management -- FAO 2020 & Organizational change management and team development for humanitarian staff. Experiential learning approach to improve team dynamics and leadership effectiveness. & Kenya (Nairobi, Turkana, Garissa) | 2020 & 120+ FAO staff trained in change management and emotional intelligence. Post-training assessment showed 85\% improvement in team cohesion and 70\% improvement in change readiness scores. Reference: FAO Kenya HR Director. \\
\hline
3 & Sustainability of Business Grants Given to Pastoralist Women -- Boma Project 2018 & Business development and sustainability training for pastoralist women entrepreneurs. Market access support and mentorship for women-led micro-enterprises. & Kenya (Turkana, Marsabit) | 2018 & 150+ pastoralist women trained in business sustainability. 60\% of participants reported increased revenue within 6 months. Reference: Boma Project Programme Manager. \\
\hline
4 & Inspiring the Next Generation of Young Africans -- Young Africa Leadership Initiative (YALI) 2017 & Leadership development and entrepreneurship training for young African leaders. Platform for peer learning, mentorship, and cross-border collaboration. & East Africa (Kenya, Uganda, Tanzania) | 2017 & 200+ young leaders trained in leadership, entrepreneurship, and digital skills. 40\% launched or scaled social enterprises within 12 months. Reference: YALI East Africa Regional Director. \\
\hline
5 & SANARA Program -- Mastercard Foundation / HEVA Fund & Youth empowerment through creative and cultural industries. Platform-mediated marketplace for creative services, business incubation, and employment readiness. & Kenya (Nairobi, Turkana, Garissa) | 2022--2025 & 600+ youth trained in creative and digital skills. 45\% women, 40\% refugees, 5\% PWDs. Platform pilot demonstrated credible marketplace demand with 200+ active providers and 500+ client engagements. Reference: HEVA Fund Programme Director. \\
\end{longtable}

% SECTION 9: M&E
\section{Monitoring, Evaluation, and Reporting}
\begin{longtable}{@{}p{4cm}p{4cm}p{4cm}p{3cm}@{}}
\textbf{Outcome to Track} & \textbf{Indicator} & \textbf{Verification Method} & \textbf{Reporting Frequency} \\
\hline
Participants assessed & Number of candidates assessed against platform readiness criteria & Platform registration data, assessment scores, facilitator reports & Monthly \\
\hline
Participants onboarded & Number of providers with complete profiles on platform & Platform user database, profile completion audit & Monthly \\
\hline
Participants activated & Number of providers who completed first paid engagement & Platform transaction records, client confirmation, payment receipts & Monthly \\
\hline
Retention at Month 1 & \% of activated providers still active after 30 days & Platform activity logs, earnings data & Monthly \\
\hline
Retention at Month 3 & \% of activated providers still active after 90 days & Platform activity logs, earnings data, facilitator check-in records & Quarterly \\
\hline
Retention at Month 6 & \% of activated providers still active after 180 days & Platform activity logs, earnings data, facilitator check-in records & Quarterly \\
\hline
Earnings outcomes & Average monthly earnings per active provider & Platform payment records, provider self-reporting & Quarterly \\
\hline
Client satisfaction & Average client rating per provider, repeat client rate & Platform review system, client survey & Quarterly \\
\hline
Inclusion outcomes & \% women, refugees, PWDs among activated providers & Platform demographic data, registration records & Quarterly \\
\hline
Platform demand & Number of active clients, booking volume, repeat usage & Platform client database, transaction records, booking analytics & Monthly \\
\end{longtable}

% SECTION 10: DECLARATION
\section{Declaration and Authorised Signature}
\noindent By submitting this proposal, the undersigned confirms that the information provided is accurate and complete, that the organisation meets all eligibility requirements set out in the Terms of Reference, and that the organisation commits to the obligations set out therein including ILO's standards on ethics, gender equality, non-discrimination, data protection, and safeguarding.

\vspace{1cm}

\begin{tabular}{ll}
\textbf{Organisation Name:} & INK Coaching Limited \\
\textbf{Authorised Signatory:} & Patrick Njiru Nyaga, Lead Consultant \\
\textbf{Signature:} & [Signature to be provided] \\
\textbf{Date:} & 10 July 2026 \\
\end{tabular}

\newpage

% FINANCIAL PROPOSAL
\section{Financial Proposal}
\noindent \textbf{Total Budget Request: USD 265,625}

\subsection{Budget Breakdown by Deliverable}
\begin{longtable}{@{}lrr@{}}
\textbf{Deliverable} & \textbf{Amount (USD)} & \textbf{\% of Total} \\
\hline
D1: Inception report and validated demand & 79,687 & 30\% \\
D2: Mid-term progress report and $\geq$50\% activation & 106,250 & 40\% \\
D3: Monthly activation progress reports (Months 4--10) & 26,563 & 10\% \\
D4: Quarterly retention and performance reports & 26,563 & 10\% \\
D5: Final report and recommendations & 26,562 & 10\% \\
\hline
\textbf{Total} & \textbf{265,625} & \textbf{100\%} \\
\end{longtable}

\subsection{Platform Commission Structure}
\noindent The VUKA platform charges a \textbf{10\% commission} on completed transactions. This fee is disclosed to all providers at registration and on every booking confirmation. Providers retain 90\% of their earnings. The commission covers platform maintenance, payment processing, dispute resolution, and provider support services.

\subsection{Company Registration Details}
\noindent [To be completed by Aunty Janet / Mom:]
\begin{itemize}[leftmargin=*]
    \item Company Registration Number: \_\_\_\_\_\_\_\_\_\_\_\_\_\_\_\_\_\_\_\_
    \item Year of Registration: \_\_\_\_\_\_\_\_\_\_\_\_\_\_\_\_\_\_\_\_
    \item Registering Body: \_\_\_\_\_\_\_\_\_\_\_\_\_\_\_\_\_\_\_\_
    \item Physical Address: \_\_\_\_\_\_\_\_\_\_\_\_\_\_\_\_\_\_\_\_
\end{itemize}

\subsection{Key Personnel Rates}
\begin{longtable}{@{}lrrr@{}}
\textbf{Role} & \textbf{Monthly Rate (USD)} & \textbf{Duration (Months)} & \textbf{Total (USD)} \\
\hline
Team Lead / Project Manager & 5,000 & 12 & 60,000 \\
Marketplace Lead / Platform Specialist & 4,000 & 12 & 48,000 \\
County Facilitator -- Turkana & 2,500 & 12 & 30,000 \\
County Facilitator -- Garissa & 2,500 & 12 & 30,000 \\
\hline
\textbf{Personnel Subtotal} & & & \textbf{168,000} \\
Platform \& Operations & & & 50,000 \\
Training \& Onboarding & & & 25,000 \\
Travel \& Logistics & & & 15,000 \\
Reporting \& M\&E & & & 7,625 \\
\hline
\textbf{Total} & & & \textbf{265,625} \\
\end{longtable}

\newpage

% BACKGROUND TEXT FROM AUNT
\section*{Background: INK Coaching and the VUKA Youth Connect Program}
\addcontentsline{toc}{section}{Background: INK Coaching and the VUKA Youth Connect Program}

\noindent At INK Coaching, we have had the privilege of participating in the implementation of the SANARA Program, a major youth empowerment initiative in Kenya that connects the creative and cultural industries with sustainable business enterprise. Inspired by the Swahili words \emph{sanaa} (arts) and \emph{biashara} (business), SANARA is backed by the Mastercard Foundation and implemented by a consortium of organizations including the Sanara -- HEVA Fund. The program places a strong emphasis on inclusion, specifically targeting women, refugees, and persons with disabilities in both urban and rural communities. Through our role within this ecosystem, INK Coaching has contributed to youth training, professional development, and employment readiness support aimed at strengthening pathways to sustainable livelihoods and work opportunities.

\noindent Building on the lessons, partnerships, and impact achieved through SANARA, we have developed an initiative called the VUKA Youth Connect Program. The proposed program is envisioned as a 28-month initiative designed to create a sustainable and inclusive ecosystem that empowers marginalized youth in Kenya to participate competitively in the digital economy.

\noindent Our goal is to bridge the gap between emerging African tech talent and the growing demand for digital skills in European markets, while promoting inclusive economic growth and cross-border innovation.

\noindent The program will specifically target youth between the ages of 18--35, including Persons with Disabilities (PWDs) and refugees living in marginalized areas across Kenya.

\noindent The VUKA Youth Connect Program will be implemented in four key phases:

\subsection{Phase 1 -- Talent Identification}
\noindent This phase will focus on identifying high-potential youth from underserved and marginalized communities through community outreach, partnerships, and localized engagement initiatives.

\subsection{Phase 2 -- Digital Skills Enhancement}
\noindent Selected participants will undergo intensive training aimed at strengthening digital competencies, remote work preparedness, employability skills, communication, and professional development aligned with international market demands.

\subsection{Phase 3 -- Employer Matching \& Entrepreneurship Incubation}
\noindent Participants will receive remote work readiness support and be connected to potential employers and virtual work opportunities. In addition, youth interested in entrepreneurship will undergo incubation support, mentorship, and business development training, culminating in an opportunity to pitch their ventures for potential funding and investment consideration.

\subsection{Phase 4 -- Sustainability Through the VUKA Platform}
\noindent To ensure long-term impact and ecosystem sustainability, we intend to establish the VUKA Platform as a continuous support and engagement hub connecting youth to opportunities, mentorship, employers, investors, trainers, and peer learning networks.

\end{document}
